
\begin{figure*}[t]
\centering
\begin{tikzpicture}[
    node distance=1.5cm,
    block/.style={rectangle, draw, rounded corners, minimum height=1cm, align=center, fill=blue!10},
    arrow/.style={thick,->,>=stealth}
]

% Input
\node (input) [block, fill=green!10] {
    \textbf{Input}\\
    Path Description\\
    (landmarks, directions)
};

% LLM
\node (llm) [block, above of=input, minimum width=4cm, fill=yellow!20] {
    \textbf{LLM (Qwen-3.5)}\\
    Prompt-based Generation
};

% Chinese Instruction
\node (instruction) [block, above of=llm, fill=green!10] {
    \textbf{Chinese Instruction}\\
    "直走 3 米，然后左转..."
};

% VLN Model
\node (vln) [block, right of=instruction, xshift=4cm, minimum width=5cm, fill=blue!20] {
    \textbf{VLN-BERT Baseline}\\
    \small Instruction Encoder + Visual Encoder\\
    \small Cross-Modal Attention
};

% Visual Features
\node (visual) [block, below of=vln, fill=orange!10] {
    \textbf{ResNet-50}\\
    Pre-trained Features\\
    (2048-d $\rightarrow$ 256-d)
};

% Output
\node (output) [block, right of=vln, xshift=3cm, fill=green!10] {
    \textbf{Action Prediction}\\
    Softmax over 36 candidates
};

% Arrows
\draw [arrow] (input) -- node[right] {Path Info} (llm);
\draw [arrow] (llm) -- node[right] {Generate} (instruction);
\draw [arrow] (instruction) -- node[above] {Tokenize} (vln);
\draw [arrow] (visual) -- node[right] {Features} (vln);
\draw [arrow] (vln) -- node[above] {Predict} (output);

\end{tikzpicture}
\caption{Overview of our LLM4VLM framework. The LLM generates Chinese navigation instructions from path descriptions, which are then fed into the VLN-BERT baseline model along with ResNet-50 visual features for action prediction.}
\label{fig:architecture}
\end{figure*}
